% ==============================================================================
% ПАРТИЯ B03-033: DISS-005-B1859 .. DISS-005-B1880
% Глава 3: Методика нормализации и разделение сигналов на группы
% ==============================================================================

% DISS-005-B1859
\section{Нормализация сигналов и преобразование форматов}\label{sec:ch3_normalization_and_formats}

% DISS-005-B1860
После этапов сбора, первичной обработки, деперсонализации и унификации имен сигналов, следующим важным шагом в подготовке датасета реальных осциллограмм является их нормализация и преобразование в формат, удобный для последующего анализа и, в частности, для задач машинного обучения. Нормализация позволяет привести амплитуды различных сигналов к сопоставимым величинам, что необходимо для корректной работы многих алгоритмов. Преобразование в стандартизированный табличный формат, такой как CSV, упрощает дальнейшую обработку данных и их использование в различных программных средах.

% DISS-005-B1861
В данном разделе будут рассмотрены методика определения базисных значений, разработанная для нашего датасета, процесс применения нормализации и, наконец, преобразование осциллограмм в формат CSV (этапы реализованы с использованием классов CreateNormOsc, NormOsc и RawToCSV).

% DISS-005-B1862
\subsection{Методика определения базисных значений для нормализации}\label{sec:ch3_basis_values_methodology}

% DISS-005-B1863
Не \commentnote{DISS-005-C0123}{Алексей Евдаков [2]}{Надо будет убрать эти излишние слова про «СЛОЖНАЯ задача», а то везде пихается, а хорошо бы в 1-2 местах\dots}{менее }сложной задачей при работе с большим набором реальных осциллограмм от различных устройств и энергообъектов является определение корректных базисных (номинальных) значений для каждого сигнала. Эти значения необходимы для приведения амплитуд к относительным единицам, что позволяет сравнивать сигналы между собой и корректно применять пороговые алгоритмы. Указанный шаг крайне важен для дальнейшего использования для задач машинного обучения. Данная методика подробно описана далее (реализованная в классе CreateNormOsc)

% DISS-005-B1864
\subsubsection{Разделение сигналов на группы}\label{sec:ch3_signal_groups}

% DISS-005-B1865
На основе стандартизированных имен (раздел~3.4.4.2) аналоговые сигналы делятся на логические группы, так как для каждой группы могут применяться свои номинальные значения и логика их определения делится на следующие группы, которые далее и обрабатываются:

\begin{itemize}
% DISS-005-B1866
\item Напряжения на сборных шинах (например, U~|~BusBar-1~|~phase:~A).
% DISS-005-B1867
\item {\emergencystretch=2em Напряжения на кабельных/воздушных линиях (например, U~|~CableLine-1~|~phase:~A).\par}
% DISS-005-B1868
\item Фазные токи (например, I~|~Bus-1~|~phase:~A).
% DISS-005-B1869
\item Ток нулевой последовательности (например, I~|~Bus-1~|~phase:~N).
% DISS-005-B1870
\item Токи дифференциальных защит (например, I~|~dif-1~|~phase:~A).
\end{itemize}

% DISS-005-B1871
\noindent\textbf{\textit{Примечание}}: \textit{сигналы дифференциальных защит часто требуют особого подхода к нормализации или могут не нормализоваться вовсе, в зависимости от конкретной задачи анализа.}

% DISS-005-B1872
\subsubsection{Использование первичных и вторичных значений}\label{sec:ch3_primary_secondary_values}

% DISS-005-B1873
Стандартными вторичными значениями для измерительных трансформаторов в странах СНГ являются 100~В для трансформаторов напряжения (ТН) и 1~А или 5~А для трансформаторов тока (ТТ). В нашем датасете для ТТ по умолчанию принимается базис 5~А, так как это наиболее распространенное значение, и точно определить его из осциллограммы часто невозможно. В последующем эти значения в полу-ручном режиме были уточнены.

% DISS-005-B1874
Коэффициенты трансформации (primary и secondary из .cfg файла) используются для пересчета между первичными и вторичными величинами, если это необходимо и если эти коэффициенты корректно указаны.

% DISS-005-B1875
\subsubsection{Алгоритм определения номинала (базиса) для каждого типа сигнала}\label{sec:ch3_basis_determination_algorithm}

% DISS-005-B1876
Изначально предполагалось, что наличие суффикса «\_raw» в имени сигнала (например, U\_raw~|~BusBar-1~|~phase:~A) указывает на использование нетрадиционных датчиков (катушки Роговского для тока, емкостные/резистивные делители для напряжения). Такие сигналы могут требовать специфической логики обработки (например, интегрирования для катушек Роговского). Однако в текущей реализации датасета все сигналы, включая, прошли через общий процесс определения базиса и нормализации, но сам флаг сохранен для возможности будущего раздельного анализа.

% DISS-005-B1877
{\emergencystretch=1.5em Для каждого аналогового сигнала выполняется быстрое преобразование Фурье (БПФ) на скользящем окне (обычно равном одному периоду промышленной частоты, например, 32~точки при частоте дискретизации 1600~Гц для сети 50~Гц). Рассчитывается амплитуда первой гармоники ($H_1$) и максимальная амплитуда среди всех остальных значимых гармоник ($H_{x\text{\_}\max}$, например, со 2-й по 15-ю).\par}

% DISS-005-B1878
На основе абсолютных значений $H_1$ (во вторичных величинах, если возможно) и соотношения $H_1/\allowbreak H_{x\text{\_}\max}$ принимается решение о базисном значении. Логика этих алгоритмов приведена в таблицах~3--5.

% DISS-005-B1879
Отметим, что данный алгоритм по напряжению в теории может привести к некорректному определению номиналу при однофазных замыканиях на землю (ОЗЗ). Так как 140~В, это перенапряжение равное 2,43 крата (140/57,73). Поэтому более подробно было осуществлено исследование всего датасета на наличие ОЗЗ и анализ осциллограмм, у которых было установлен номинал 400~В. Текущий анализ показал, что таких случаев перенапряжений не имеется, что говорит о крайне низкой вероятности их возникновения (более подробно этот вопрос рассмотрен в\commentnote{DISS-005-C0124}{Алексей Евдаков [2]}{Как раз будет проведено отдельное исследование и будет показана актуальность данного набора не только для нейронок}{} главе ХХ)

% DISS-005-B1880
\paragraph{Таблица 3. \commentnote{DISS-005-C0125}{Алексей Евдаков [2]}{Хорошая идея от нейронок – заменить на блок схемы. Ведь в коде как раз так и прописано. Это надо перерисовывать, поэтому пока лишь озвучу как идею}{Алгоритм }автоматического определения базисного значения напряжения}\label{tab:ch03_voltage_basis_algorithm}

% ==============================================================================
% ПАРТИЯ B03-034: DISS-005-B1881 .. DISS-005-B1898
% Глава 3: Таблица 3 (Алгоритм автоматического определения базисного значения напряжения, TAB-083, 18 ячеек)
% ==============================================================================

\mbox{}\par\vspace{0.2cm}
{\small
\centering
\renewcommand{\arraystretch}{1.3}
\noindent
\begin{tabular}{|>{\centering\arraybackslash}p{105pt}|>{\centering\arraybackslash}p{95pt}|>{\centering\arraybackslash}p{95pt}|>{\centering\arraybackslash}p{115pt}|}
\hline
% DISS-005-B1881
\multirow{2}{*}{Условия} &
% DISS-005-B1882
\multicolumn{2}{c|}{$H_{1(s)} > 1{,}5\,H_{x\text{-max}(s)}$} &
% DISS-005-B1883
\multirow{2}{*}{$H_{1(s)} \le 1{,}5\,H_{x\text{-max}(s)}$} \\ \cline{2-3}
% DISS-005-B1884
&
% DISS-005-B1885
$H_{1(p)} > 560$ &
% DISS-005-B1886
$H_{1(p)} \le 560$ &
% DISS-005-B1887
\\ \hline
% DISS-005-B1888
$H_{1(s)} > 560$ &
% DISS-005-B1889
\multicolumn{3}{c|}{?3} \\ \hline
% DISS-005-B1890
$140 < H_{1(s)} \le 560$ &
% DISS-005-B1891
\multicolumn{2}{c|}{$U_{\text{base}} = 400$} &
% DISS-005-B1892
?2 \\ \hline
% DISS-005-B1893
$20 < H_{1(s)} \le 140$ &
% DISS-005-B1894
\multicolumn{2}{c|}{$U_{\text{base}} = 100$} &
% DISS-005-B1895
?2 \\ \hline
% DISS-005-B1896
$H_{1(s)} \le 20$ &
% DISS-005-B1897
?1 &
% DISS-005-B1898
\multicolumn{2}{c|}{Noise} \\ \hline
\end{tabular}
\par}
\vspace{0.3cm}

